\scanpage{29}
avec un poids marqué $\unsure{\nu}$, et bien sûr, un genre donné). Du pt de vue \unsure{structure} de graphe associé au type numérique $g,\nu)$, c'est le plus simple de tous ! On étudiera les plus complexes tantôt, qui correspondent au cas d'une dimension modulaire nulle, \unsure{ou} d'une codimension modulaire maximale, égale à $3(g-1)+\nu$.

Il est tentant d'introduire le compactifié
\[
\widehat{\mathcal M}_G\overset{\mathrm{déf}}{=}\prod_{\alpha\in S}\widehat{\mathcal M}_{g_\alpha,\widehat I_\alpha}
\]
où les $\widehat{\mathcal M}_{g_\alpha,\widehat I_\alpha}$ sont les compactifiés de MD des $\mathcal M_{g_\alpha,\widehat I_\alpha}\simeq\mathcal M_{g_\alpha,\widehat\nu_\alpha}$. \deleted{\unsure{Ils|Il}}\inserted{Pour $\alpha$ fixé, $\widehat{\mathcal M}_{g_\alpha,\widehat I_\alpha}$} \unsure{représente|représentent} le groupoïde fibré
\[
(30)\qquad\begin{aligned}
\mathcal S\longmapsto&\ \text{courbes de MD relatives }\inserted{Y_\alpha}\text{ sur }\mathcal S,\\
&\text{de type }g_\alpha,\widehat I_\alpha\quad\text{i.e. de genre}\\
&g_\alpha,\text{ munies d'un mono}\\
&(\widehat I_\alpha)_{\mathcal S}\hookrightarrow Y_\alpha^{\text{lisse}}
\end{aligned}
\]
de sorte que $\widehat{\mathcal M}_G$ représente le groupoïde fibré
\mspage{28}
